% 本文件由 LaTeX 排版 Agent 根据有效 translation.json 自动生成。
% author=Lactantius; work_id=0704; title=论天主的化工

\chapter{论天主的化工}

% structure: p0001 source=h1 target=omitted reason=page-title-already-rendered-by-parent

一篇致其门徒德默特黎亚努斯的论文\par

% structure: p0003 source=h2 target=section reason=relative-heading-rank; base=h2

\section{第一章：引言及对德默特黎亚努斯的劝勉。}

我多么烦乱、处于极大的困境中，你将能从这本我写给你的小书中判断出来，德默特黎亚努斯，我几乎是用不加修饰的话语写的，正如我才能的平庸所允许的那样，为叫你知道我每日的追求，也让我不至于亏欠你——即使现在仍作为导师，却是关于更尊荣的主题和更好的体系。因为如果你在文学——那只不过训练文体——上甘愿作一个听话的听众，那么在这些关乎生命的真正学问中，你更应受教！我现在向你声明：没有任何环境或时间的障碍能阻止我编写一些东西，使我们所持守的学派的哲学家们将来能变得更有学识、更有教养，尽管他们如今名声不佳，常受责备，因为他们生活的方式不像智者所应为，而且以名声为掩护隐藏自己的恶习；他们本应纠正这些，或完全避免，使智慧之名因生活自身符合教训而幸福且纯洁。然而，我不躲避任何辛劳，为能同时教导自己与他人。因为我不能忘记自己，尤其在此刻，那是我最需要记住的时候；正如你也不忘记自己——我盼望并愿意如此。虽然国家的需求可能使你偏离真实和正义的工作，但一个正直的良心不可能不时常仰望上天。\par

我确实欢喜一切被认为是福乐的事都顺利归于你，但条件是它们不改变你的心境。因为我怕习惯和这些事物的快感会像通常那样逐步潜入你的心。因此我劝告你，\par

\Quote{并且重复地，将再三劝告你，}\par

不要相信你拥有这些地上的享受是伟大或真正的福乐，因为它们不仅因不确定而欺骗人，而且因愉悦而阴险。因为你知道那个摔跤手和我们的对手是多么狡猾，也时常暴力，正如我们现在看到他那样。他利用一切能诱惑的东西作为陷阱，而且如此巧妙，以致逃避了心智的眼睛，使人无法凭预见避免。因此，最高的智慧是步步为营，因为他占领了通道的两侧，并暗中为我们的脚放置绊脚石。因此我劝告你：要么，如果你能按你的德行，蔑视你生活在其中的顺境；要么，不要过分赞赏它。记住你真正的父母，记住你是在哪座城里报的名，以及你曾属于何等阶层。你肯定明白我所说的。因为我不指控你骄傲——你甚至没有一丝怀疑；但我所说的是指向心灵，而非身体——整个身体的构造为此安排，为叫它服从灵魂如同主人，并由灵魂的意志所支配。因为它在某种意义上是陶制的器皿，其中盛着灵魂——即真正的人自己——而且这器皿并非如诗人所说由普罗米修斯所造，而是由那位至高无上的创造者和世界的工匠——天主所造；天主的眷顾和最完美的卓越，既无法凭感知领悟，也无法用言语表达。\par

然而，既然提到了身体和灵魂，我将尽力按我的理解力所见的微弱光亮，解释二者的本性；我认为这项责任尤其应当承担，因为马尔库斯·图利乌斯（Marcus Tullius），一位才华卓越的人，在他的\WorkTitle{论共和国}第四卷中，曾试图做这件事，却将广阔的主题局限在狭窄的范围内，轻微地选取了要点。而为了不找借口说他没有继续这个主题，他证明自己既非缺乏意愿也非缺乏关注。因为在他的\WorkTitle{论法律}第一卷中，当简要总结同一主题时，他这样说：\Quote{西庇阿，在我看来，已经在你读过的那几本书中充分表达了这一主题。}后来，在他的\WorkTitle{论神性}第二卷中，他试图更广泛地继续同一主题。但既然他连那里也没有充分表达，我将从事这项职责，并大胆地承担解释那位最雄辩之人几乎未曾触及之事。也许你会责备我试图讨论一些晦涩之事，当你看到曾有如此鲁莽的人（通常被称为哲学家）时，他们探究了天主愿为深奥隐藏之事，考察天上地下的本性——这些离我们很远，无法用眼看见，用手触摸，用感官感知；然而他们却如此争辩这些本性，以致希望他们提出的东西看起来是证明过的和已知的。请问，有什么理由使人认为我们若愿意考察并默观我们身体的结构（这并非完全晦暗，因为从肢体的功能和各部分的使用，我们可以理解每一部分是以何等大的天主的眷顾之能造出的）就是可憎的呢？\par

% structure: p0009 source=h2 target=section reason=relative-heading-rank; base=h2

\section{第二章：论野兽与人的受造。}

因为我们的造物主和父亲，天主，赐予人感知和理性，由此显明我们出自祂，因为祂本身就是理智，祂本身就是感知和理性。既然祂未将那种理性的能力赐给其他动物，便预先安排它们的生活方式更为安全。祂让它们全都披上自身的天然毛发，以便更容易忍受霜雪与寒冷的严酷。此外，祂给每一物种指定了独特的防御，以抵御外来的攻击：它们或能以自然武器对抗更强的动物，或能以迅捷的逃跑脱离危险，或需兼具力量与速度者能以巧计自保，或在藏身之处守卫自己。因此，它们中有的以轻羽升腾，有的以蹄支撑，有的以角武装；有的口中有利器——即牙齿——或脚上有钩爪；没有哪一个缺少自卫的防御。\par

倘若有动物沦为更大动物的猎物，为使它们不致完全灭绝，它们或被放逐到那些较大动物无法生存的区域，或得以享有更丰盛的繁殖力，以便为靠血滋养的野兽提供食物，同时它们自身的众多数量仍能在遭受的杀戮中存活，以保存物种。但祂造人时——将理性赐予他，将感知和言说的能力赋予他——却使他缺乏赐给其他动物的那些东西，因为智慧能够弥补本性条件所拒绝给他的东西。祂使他赤裸且无防卫，以便他可以凭才能武装自己，凭理性遮盖自己。然而，那些赐给野兽的东西的缺失，竟如此奇妙地有助于人的优美，这无法言表。因为如果祂赐给人野兽的牙齿、角、爪、蹄、尾巴或各色毛发，谁能看不出这是一个何等畸形的动物——正如哑巴动物若是赤裸无防卫，也会如此丑陋？因为若从它们身上取去其体表的天然衣裳，或那些用于自身武装的东西，它们既不能美观也不能安全；以至于它们若从功用方面看显得奇妙地装备，从外观方面看又显得奇妙地装饰：功用与美丽如此奇妙地结合。\par

至于人，祂将他造成一个永恒和不朽的受造物，并未如其他动物那样从外面装备他，而是从里面；祂将他的保障不放在身体中，而放在灵魂里；因为既然祂已赐给他那至高的价值，再以身体上的防御遮盖他便是多余的——尤其当这些防御妨害人体之美时。为此，我常对追随厄壁鸠鲁的哲学家的愚昧感到惊异；他们责难自然的作品，试图表明世界并非由任何上智安排所预备和治理，却将万物的起源归于不可分割的坚固物体，说万物皆由这些物体的偶然碰撞而产生和生成。我暂且不提他们挑剔的作品本身——那方面他们荒谬得疯狂；我只涉及与我们当前主题直接相关的内容。\par

% structure: p0013 source=h2 target=section reason=relative-heading-rank; base=h2

\section{第三章——论野兽与人的境况}

他们抱怨说，人生来比其它动物的出生更软弱、更脆弱：因为其它动物一经从子宫产出，立刻用脚站起，以来回奔跑表达喜悦，并立即能承受空气——它们既已带着天生的遮盖来到光中；而人却相反，赤身无防卫，仿佛从船难中被抛到今生之苦海中；他既不能从出生处移动自己，也不能寻求乳汁的营养，更不能忍受时日的侵害。因此他们说，自然不是人类之母，而是继母，她如此慷慨地对待哑巴动物，却这样生出人：毫无资源、毫无力量、丧失一切援助，他除了以哭号和哀叹表达其脆弱状态外，别无他法；\Quote{他所要经历的，正是他命定要在生命中承受如此多苦难的命运。}\par

当他们说这些话时，人们相信他们非常明智，因为每个人不加思索便对自己的处境不满；但我认为，他们从未像说这些话时那样愚蠢。因为当我思量万物的状况时，我明白没有什么是本该不同的——不是说能够不同，因为天主无所不能；而是必须如此：那最富有上智安排的尊严者，必定创造了更好、更正确的事物。\par

因此，我愿意询问那些指责天主的作品的人，他们认为人身上因出生时更软弱的状况而缺少什么？他们认为人因此被抚养得更差吗？或者他们较少达到年龄的最大力量？还是软弱阻碍了他们的成长或安全，因为理性赐予了所欠缺的东西？但是，他们说，抚养人需要最大的劳苦：事实上，野兽的条件更好，因为这些野兽在生育幼崽后，除了自己的食物外毫无挂虑；由此，它们的乳房自动膨胀，乳汁的滋养供应给后代，而幼崽出于本性的强迫寻找这滋养，母亲毫无劳苦。飞鸟的情况如何，它们有不同本性？它们在抚养幼鸟时不也承受最大的劳苦，以至于有时似乎拥有某种人类的智慧？因为它们或用泥土筑巢，或用树枝和树叶建造，并且不吃食物地伏在蛋上；既然没有赋予它们用自己的身体养育幼鸟的能力，它们就搬运食物给幼鸟，整天这样来回奔波；但夜间它们保卫、爱护、保护它们。人还能做什么更多的事？除非只是这一件，就是他们不赶走长大的幼崽，而是通过永久的关系和情感的纽带留住他们。我何必说飞鸟的后代比人的后代脆弱得多呢？因为它们不是从母体中生出动物本身，而是那通过母体身体的滋养和热量得到温暖而产生动物的东西；并且这即使有了气息，由于无羽和柔弱，不仅不能飞，甚至不能行走。因此，如果有人以为自然对待飞鸟很糟糕，首先因为它们两次出生，然后因为它们如此软弱，以至于必须由父母辛苦觅食来喂养它们，难道他不是最无理智的吗？但是，他们选择较强的，而忽略较弱的动物。\par

因此，我问那些偏好野兽状况胜过自身状况的人，如果天主给予他们选择，他们会选择什么：是选择有软弱的智慧，还是选择有本性的力量？事实上，他们并不如此像野兽，以至于不愿选择甚至更软弱的状况，只要那是人的，而不愿选择那没有理性的力量。但是，实话说，审慎的人既不想要有软弱的理性，也不想要没有理性的哑巴动物的力量。因此，没有如此抵触或矛盾之事，以至于理性或本性的状况必然为每种动物准备一切。如果它配备有自然的保护，理性就是多余的。因为它会发明什么？会做什么？会计划什么？或者会在哪里显示那理智之光，当本性自动就赐予那些能够是理性结果的东西时？但如果它被赋予理性，身体的防御又有什么需要呢？因为一旦赋予理性，它就能够代替本性的功能。而且，这理性对于人的装饰和保护有如此大的力量，以至于天主不能赐予更大或更好的东西。最后，由于人的身体不大、力量不强、健康脆弱，然而，因为他接受了更有价值的东西，他就比其他动物装备得更好，装饰得更美。因为尽管他生来脆弱软弱，但却安全地躲避所有哑巴动物，而那些出生时更有力量的动物，虽然能够忍耐气候的严酷，却无法安全地躲避人。因此，理性赐予人的多于自然赐予哑巴动物的；因为，在它们的情况下，无论力量的大小还是身体的坚固，都不能阻止它们被我们压迫，或被迫服从于我们的权力。\par

那么，任何人当他看到甚至大象以其巨大的身体和力量都服从于人时，还能抱怨造物主天主，因为他只接受了适中的力量和矮小的身体，而不按照其应得的评价天主的恩惠吗？这正是不感恩的人的部分，或更真实地说，是疯子的部分。\Person{柏拉图}，我相信，为了反驳这些不感恩的人，感谢自然因为他生为人。他做得多么更好更正确，他认识到人的状况比那些宁愿自己生为野兽的人更好！因为如果天主碰巧把他们变成那些他们偏好其状况甚于自身的动物，他们现在会立刻希望回到以前的状态，并且会大声呼喊急切地要求他们以前的状况，因为身体的力量和坚固并不重要到让你没有舌头的作用；或者飞鸟在空中的自由飞行并不重要到让你没有双手。因为双手比翅膀的轻捷和用处更有益处；舌头比整个身体的力量更有益处。因此，偏好那些如果被赐予你就会拒绝接受的东西，是何等疯狂！\par

% structure: p0019 source=h2 target=section reason=relative-heading-rank; base=h2

\section{第四章 论人的软弱}

他们又抱怨人易于患病和夭折。他们显然愤慨自己不是生而为神。他们说：\Quote{绝非如此；而是我们由此表明，人的受造毫无远见，本应不同。}若我指出，正是这件事被极有理由地预见到，即他可能受疾病折磨，且其生命常在途中被截短，那又如何？因为天主既已知道祂所造的生物自行趋向死亡，为使它能承受死亡本身——即本性的解体——便赐予它脆弱性，使死亡有一入口，以达成生物的解体。因为若它如此强壮，以致疾病和病痛无法接近，那么死亡也不能接近，因为死亡是疾病的结果。但对于那已被指定了成熟死亡的人，夭折又如何能远离他呢？无疑，他们希望无人死亡，除非他已活满一百岁。在如此对立的情况下，他们怎能保持一致性？因为，为使无人能在百岁前死亡，必须赐予他某种不朽的力量；而一旦赐予，死亡的条件必然被排除。但那能使人坚固、抵御疾病和外来攻击的力量，会是怎样的呢？因为，既然人由骨骼、神经、肉和血组成，其中哪一样能如此坚固以抵御脆弱和死亡？因此，为使人在他们本应为他指定的时间之前不致解体，他们要用什么材料来赋予他一个身体？所有能看见和触摸的事物都是脆弱的。他们只得从天上寻求某物，因为地上没有不脆弱的东西。\par

既然人不得不被天主如此塑造，即他终有一死，事情本身要求他应以脆弱和属地的身体被造。因此，他终有一死是必然的，因为他拥有一个身体；因为每个身体都可解体并死亡。因此，那些抱怨夭折的人极其愚蠢，因为本性的状况为它留出了余地。如此，他也会受疾病支配；因为本性不允许那终将解体的身体没有软弱。但让我们假设，如他们所愿，人并非生来就处于受疾病或死亡支配的条件下，除非他走完生命历程，到达年迈的尽头。因此，他们没有看到，如果作此安排，其后果将是什么：显然不可能在别的时间死亡；但如果有人能被他人剥夺营养，他就有可能死亡。因此，情况要求人——不能在指定日期之前死亡——不应需要食物的滋养，因为食物可能被夺走；但若他不再需要食物，他将不再是人了，而成为神。因此，如我已说的，那些抱怨人的脆弱的人，尤其抱怨他们不是生而不朽和永存的。除非年老，无人应当死亡。正因如此，他确实应当死亡，因为他不是天主。但死亡不能与不朽结合：因为如果人在老年是可死的，他就不能在青年时不朽；死亡的条件对于终将死去的人并不陌生；也没有任何不朽被设定了限度。因此，永远的排除不朽和暂时的接受死亡，使人处于这样一种状况：他终有一死。\par

因此，在所有方面，必然性都是恰当的：他不应不同于他现在的样子，而且这也是不可能的。但他们看不到后果的顺序，因为他们已在大问题上犯了错误。因为天主的眷顾被排除在人事之外，就必然导致万物自行产生。因此，他们发明了那种微小种子的碰撞和偶然相遇的概念，因为他们看不到事物的起源。当他们陷入这一困境时，必然性迫使他们认为灵魂与身体一同出生，也同样与身体一同消灭；因为他们假设了没有东西是由神圣心智所造的。他们无法用别的方式证明这一点，除非表明有些事上眷顾的系统似乎有缺陷。因此，他们责备那些眷顾奇妙地表达其神性的事，如我所述关于疾病和夭折的事；然而他们本应考虑，假设这些事，什么是必然的后果（但我所说的是后果）——如果他不受疾病支配，也不需要住所和衣物。因为他为什么要害怕风、雨、寒冷，其力量在于带来疾病？正因如此，他领受了智慧，以保护他的脆弱免受伤害之物。必然的后果是，既然他为了保持智慧而容易患病，他也必须易于死亡；因为死亡不临到的人，必然是坚固的。但脆弱本身具有死亡的条件；而在坚固之处，既没有老年的位置，也没有随之而来的死亡。\par

此外，若死亡被指定于一个固定的年龄，人就会变得极其傲慢，并丧失所有的人性。因为几乎所有使人彼此联合的人性权利，都源于恐惧和对脆弱性的认识。简言之，所有较软弱且胆怯的动物都聚集在一起，以便它们由于不能以力量保护自己，就可以靠数量来自卫；但较强的动物则寻求独处，因为它们信赖自己的力量和勇力。倘若人也同样有足够的力量抵御危险，不需要任何他人的帮助，那还有什么社会？或什么制度？什么人性？或有什么比人更冷酷？更残忍？更野蛮？但正因他软弱，不能离群索居，他渴望社会，使他与别人交往的生活变得更体面、也更安全。因此，你看，人的全部理性最集中地在于这一点：他生来赤身露体、脆弱，受疾病攻击，被夭折惩罚。如果这些从人身上被拿走，那么理性和智慧也必然被拿走。但我对这些显而易见的事讨论得太久了，因为显然没有任何事物曾经或能够在没有天主眷顾的情况下被造。如果我现在想按顺序讨论它所有的作品，这主题将是无尽的。但我只打算就人的身体谈论这么多，以便在其中显示天主眷顾的能力，仅在那些易于理解和显而易见的事物上何等伟大；因为与灵魂相关的事物既不能为眼睛所看见，也不能被理解。现在我们讨论我们所见的人本身的器皿。\par

% structure: p0024 source=h2 target=section reason=relative-heading-rank; base=h2

\section{第五章——论动物的形状和肢体。}

起初，当天主塑造动物时，祂不愿将它们聚集成球形，以便它们能够容易地移动行走和向任何方向转动；而是从身体的最上部延长出头部。祂还将一些肢体延长得更长，这些肢体被称为脚，以便它们通过交替运动固定在地上，引领动物去往其倾向或寻觅食物的需要所召唤之处。此外，祂从身体本身的器皿上使四个肢体伸出：两个在后面，这是所有动物都有的——脚；还有两个靠近头和颈，为动物提供各种用途。因为在牛和野兽中，它们像后脚一样是脚；但在人身上，它们是手，生出不是为了行走，而是为了行动和主宰。还有第三类，其中那些前肢既不是脚也不是手，而是翅膀，它们具有排列整齐的羽毛，提供飞行的用途。如此，一个形状有不同的形式和用途；并且为了牢固地保持身体本身的密度，通过将较大和较小的骨骼绑在一起，祂压实了一种龙骨，我们称之为脊柱；祂认为不应当用一块连续的骨头来形成它，以免动物不能行走和弯腰。从其中部，祂几乎像向不同方向延伸出横的和扁平的骨头，它们稍微弯曲，几乎像围成一个圆圈，以便覆盖内部器官，使那些需要柔软、不那么坚强的部分可以通过固体框架的包围而得到保护。在那个我们说过类似船龙骨的连结的末端，祂安置了头部，其中可能有整个生物的统治权；正如\Person{瓦罗}写信给\Person{西塞罗}所说，这个名字被赋予它，因为感觉和神经由此开始。\par

然而，我们说过，那些从身体延伸出来的部分，或为行走，或为行动，或为飞翔，祂使它们由骨骼构成，既不宜过长（以求动作敏捷），也不宜过短（以求稳固），而是为数不多且块头较大。它们要么如人一般有两肢，要么如四足动物一般有四肢。祂并没有把这些骨骼做成实心的，以免行走时迟缓沉重而拖累；而是将其做成中空，内里充满骨髓，以保持身体的活力。再者，祂并没有将它们等长地延伸到末端，而是在它们的末端聚拢成粗大的结节，以便更容易与肌腱连接，也更易于转动，由此这些部分被称为关节。这些结节，祂造得坚固实心，并覆盖一层柔软的包裹物，称为软骨；目的是为了使它们弯曲时不会擦伤或引起任何痛感。然而，祂并未按同一式样造出所有这些关节。祂使其中一些单纯而圆如球状，至少在那些肢体需要向各个方向转动的关节处，例如肩关节，因为手需要向任何方向移动和扭转；但另一些关节，祂造得宽阔、平整，并向一侧圆滑，显然是在那些只需肢体弯曲的地方，例如膝、肘和手本身。因为，手从其生长处向各方向运动，既悦目又有用；但若肘部也是如此，那样的动作便立即显得多余且不雅。因为那时，手会因过度灵活而失去现今的尊严，看起来像大象的鼻子；人就会完全变成蛇手——这一情形在那巨兽身上奇妙地实现了。因为天主愿意通过万物的奇妙多样来显示祂的眷顾和大能，既没有将那头动物的头部延伸得很长，使它能够用嘴触地（那将是可怕而丑陋的），又因为祂用伸长的獠牙武装了那张嘴，即使它触地，獠牙也会使它无法进食，于是祂在獠牙之间从前额顶部延伸出一条柔软而灵活的肢体，使它能抓取任何东西，以免獠牙的突出巨大或颈部的短浅妨碍进食的安排。\par

% structure: p0027 source=h2 target=section reason=relative-heading-rank; base=h2

\section{第六章——论\Person{伊壁鸠鲁}的错误，以及肢体及其用途}

我在此无法阻止自己再次指出\Person{伊壁鸠鲁}的愚妄。因为\Person{卢克莱修}的一切狂言都属于他；\Person{卢克莱修}为了表明动物并非出于神圣心智的巧计（如他常说的那样，而是出于偶然），便说在世界的起初产生了无数其他形态奇妙、大小各异的动物，但它们无法持久，因为要么获取食物的能力，要么交配和生殖的方式有所欠缺。显然，为了给他那在无垠虚空中飞窜的原子留出地盘，他想要排除天主的眷顾。但当他看到一种奇妙的眷顾体系存在于一切有气息之物中时，他却说曾有过体型巨大的动物，其繁衍体系已经终止——这是何等虚妄（哦，这为害之人）！\par

因此，既然我们所见到的一切都是按照一个计划产生的——因为唯有计划才能实现这种出生的条件——那么显然，没有计划就什么也不能出生。因为在形成万物的过程中，早已预见了它应如何利用肢体的服务来满足生活所需；以及后代如何通过身体的结合而产生，从而以其各自的物种保存所有生物。因为如果一位巧匠在计划建造某座大厦时，首先考虑的是整座建筑完工后的效果，并通过测量预先确定轻者适合什么位置，结构中的沉重部分应立在何处，柱子之间应留多少间隔，落水的下泄口和蓄水池应在何处或如何设置——我说，他首先预见到这些事，以便在打地基的同时就开始准备工程完成后所需的一切——那么，为什么有人会认为，在生物的构造中，天主在赋予生命本身之前，没有预见到生命所需的事物呢？因为很明显，除非那些赖以存在的事物预先安排妥当，生命就无法存在。\par

因此，\Person{伊壁鸠鲁}在动物的身体里看到了神圣计划的巧艺；但为了贯彻他先前愚昧地假定的事，他又添加了另一个与前者一致的荒谬说法。他说，眼睛不是为了看而被造的，耳朵不是为了听，脚不是为了走，因为这些器官在行使看、听、走的功能之前就已经存在；但所有这些器官的功能是在它们被造之后才从它们身上产生的。我担心，驳斥如此荒唐可笑的故事本身似乎也是同样愚蠢；但我乐于当个傻瓜，因为我们正在与一个愚昧的人打交道，免得他自以为太聪明。你怎么说，\Person{伊壁鸠鲁}？眼睛不就是为了看而被造的吗？那么，它们为何能看见呢？他说，它们的用途后来才显现出来。因此，它们就是为了看而被造的，因为它们除了看什么也做不了。同样，对于其他肢体，功能本身表明了它们是为了什么目的而被造的。因为显而易见，除非所有肢体都如此布置和预见地造成，以便能够具有其功能，否则这种功能根本不可能存在。\par

因为倘若你说，飞鸟并非为飞翔而造，野兽并非为凶猛而造，鱼类并非为游动而造，人类并非为智慧而造，那么既然显然生物都受制于各自被造时所赋予的本性及职分，这又如何解释呢？但显然，凡已失去真理本身要旨的人，必常陷于错误。因为如果万物不是由天主眷顾所生，而是由原子的偶然碰撞所产生，那么为何从未偶然发生这样的事：那些第一原理如此结合，以致造出一种能用鼻孔听、用眼睛嗅、用耳朵看的动物呢？因为如果第一原理不放过任何位置组合，那么这类怪异的产物本应每日产生，其中肢体的排列会扭曲，用途也会与通常情况大相径庭。然而，既然一切动物种类及一切肢体都遵守自身的法则、排列及所赋予的用途，显然没有什么是偶然造成的，因为天主计划的永恒秩序得以保存。但我们将在另一时机驳斥伊壁鸠鲁。现在，让我们按已开始的那样讨论天主眷顾的主题。\par

% structure: p0032 source=h2 target=section reason=relative-heading-rank; base=h2

\section{第七章——论身体的一切部分}

因此，天主连结并束缚了那些强壮身体的部分，我们称之为骨骼，它们由筋腱彼此打结并连接，这样灵魂便可如同使用绳索一般，在想要前行或迟延时加以利用；而且确实，无需任何劳作或努力，只凭极轻微的倾向，便能调节并引导整个身体的质量。祂用适合各部位的内脏覆盖它们，使坚实的部分得以隐藏起来。祂还将血管如同溪流般分布全身，与内脏混合，使湿气和血液向不同方向流动，以生命的汁液滋润所有肢体；祂按每种类型和位置所适宜的方式塑造了这些内脏，并以覆盖其上的皮肤包裹它们，这皮肤或仅以美丽为装饰，或以厚毛覆盖，或以鳞片为甲，或以绚丽的羽毛点缀。但天主奇妙的创造在于：同一种结构与状态竟显现出无穷多样的动物。因为几乎所有有气息之物中，肢体的连接与排列都是相同的。首先是头部，接着是颈部；颈部之后是胸部，胸部两侧伸出肩膀，腹部紧连胸部；腹部之下是生殖器官；最后是大腿和脚。不仅肢体在一切动物中保持自身的顺序和位置，而且肢体的各部分也是如此。因为仅就头部而言，耳朵占有固定位置，眼睛占有固定位置，鼻孔亦然，口及牙齿舌头亦然。尽管这些在一切动物中都相同，但被造物却有无限多样的差异；因为我所提及的那些部分，或更伸展或更收缩，以各种不同的线条被容纳。不仅如此，难道这不神圣吗？在如此众多的生物中，每个动物在其同类和物种中都是最优越的——因此，若将某一部分从此处移至彼处，必然的结果是：于用途而言，没有比这更不便的；于视观而言，没有比这更丑陋的；例如，若你给大象一个长颈，或给骆驼一个短颈；或者你若给蛇附上脚或毛发，而蛇那均匀伸展的身体长度本不需要别的，只需在背部带有斑点，依靠光滑的鳞片支撑，以蜿蜒的路线滑入光滑之地。然而在四足动物中，同一位设计者延长了脊柱的排列，脊柱从头顶向外延伸至身体外部更长，并收尾成尾巴，以便身体中令人不快的部分或因不雅观而被遮盖，或因柔弱而被保护，这样通过它的摆动，某些微小有害的动物就能被驱离身体；你若取下这一肢体，动物便不完美且软弱。但在有理性和手的地方，毛发的覆盖就不那么必要了。一切都如此最合宜地安排在其各自类别中，以至于无法想象比赤裸的四足动物或覆盖毛发的男子更不协调的了。\par

然而，虽然人的裸体本身在奇妙程度上趋向美丽，但它并不适合头部；因为若如此，会有多大丑陋，从秃顶便可明显看出。因此，祂用头发覆盖头部；并且因头发将在头顶，祂将其作为装饰添加，仿佛建筑的最高顶点。这装饰并非聚集成圆圈，或塑成帽子的形状，以免留下部分裸露而有碍观瞻；而是按各处的适宜性，在某些地方自由倾泻，在另一些地方收回。因此，被周围一圈环绕的前额，从太阳穴在耳朵前方生出的头发，以及这些头发的最高部分如冠冕般环绕，以及后脑全被覆盖，展现了奇妙优雅的外观。然后，胡须的性质以难以置信的方式有助于区分身体的成熟、性别的区分，或男子气概和力量之美；以致于整体工作的体系若以任何其他方式造成，便似乎不会协调。\par

% structure: p0035 source=h2 target=section reason=relative-heading-rank; base=h2

\section{第八章——论人的部分：眼与耳}

现在我要展示整个人体的构造，并解释那些在身体上显露或隐藏的各个肢体的用途和习性。因此，当天主决定在一切动物中只使人成为属天的，而其余一切属地的时，祂便使他直立以仰望上天，并造他为两足的，无疑是为使他能注视他来源的方向；但祂将其他动物压制向地，好使它们因对永生没有期望，整个身体被抛向地面，而服务于它们的食欲和食物。因此，唯独人具有的正当理性、高升的地位，以及他那与天主圣父共享并极为相似的面容，都在昭示他的起源和造主。他的心灵几乎神圣，因为它不仅获得了统治地上动物的权柄，甚至统治自己的身体，这心灵位于最高部分——头部，如同在高耸的城堡中，俯瞰并观察一切。祂造了这心灵之宫殿，不像哑巴动物那样拉伸延长，而是像圆球和球体，因为所有圆形都属于完美的设计和形状。因此，心灵和那神圣之火被其覆盖，如同拱顶；当祂用天然的覆盖物遮盖了其最高顶端时，祂同样以肢体必要的服务来装备和装饰被称为脸的前部。\par

首先，祂用凹陷的孔关闭了眼睛的球体，\Person{瓦罗}认为前额即由此钻孔而得名；祂要使这些孔不多不少恰为两个，因为没有比二更完美的数目；同样祂造了两只耳朵，这一双重性带来了令人难以置信的美，既因为每一部分都以相似性装饰，也因为来自两侧的声音更易被收集。因为耳朵本身的形状以奇妙的方式形成：因为祂不欲其开口裸露无盖，那将既不好看也不实用；因为声音可能飞越简单洞穴的狭窄空间而散开，除非开口本身将其约束，通过空腔回旋接收并阻止回声，正如那些小容器，借其应用，窄颈容器常得以灌满。\par

这些耳朵，其名称来源于饮入声音，\Person{维吉尔}说：\par

\Quote{我用这双耳朵饮入他的声音；}\par

或因为希腊人将声音本身称为\OriginalTerm{声音}{\GreekText{αὐδη}}。\par

\SourceRecord{Translated by William Fletcher.；Ante-Nicene Fathers；Vol. 7.；Edited by Alexander Roberts, James Donaldson, and A. Cleveland Coxe.；Buffalo, NY: Christian Literature Publishing Co.,；1886.；<http://www.newadvent.org/fathers/0704.htm>.}
